% !TeX root = ../教材题库_全部条目_转写工作稿.tex
% Source: XJ-B1-C05-EX-P185-Y1160
% 第5章 / 5.3 三角函数的图象与性质 / 教材 p.185 / PDF 第192页
% OCR未校对：公式与文字均可能有误；未完成数学TeX转写；禁止直接发给学生
% 原题及相邻原解截图（核对用）：../题目原图/XJ-B1-C05-EX-P185-Y1160-01.webp
% 本文件只在转写完成后作为题干使用；OCR 条目尚不是可靠题干。
\begin{bankitem}{5.3 例题 例题6}{第5章 / 5.3 三角函数的图象与性质 / 教材 p.185 / PDF 第192页}{XJ-B1-C05-EX-P185-Y1160}{OCR未校对：公式与文字均可能有误；未完成数学TeX转写；禁止直接发给学生}
\begin{ocrdraft}
[教材p.185 / PDF第192页]
\par
0河利用函数的单调性，比较下列各组数的大小：
\par
（1）tan(—3)，tan(一3．1)；
\par
（2）tan
\par
，tan
\par
一丌，且函数3，=tan在区间
\par
解（1）
\par
由于
\par
2
\par
一刁上单调递增，
\par
2
\par
因此
\par
tan(一3．1)<tan(一3）．
\par
一十丌<一<一<十丌，且函数3一t“，在区间（
\par
7丌7丌丌
\par
一十刁上
\par
（2）由于一
\par
十丌，
\par
单调递增，
\par
因此
\par
tan\ocrglyph{00AF}<tan
\end{ocrdraft}
\end{bankitem}
